\documentclass[11pt]{article}
\usepackage{amsmath,amssymb,amsthm}
\newtheorem{lemma}{Lemma}
\newtheorem{definition}{Definition}
\newtheorem{theorem}{Theorem}
\begin{document}

\begin{lemma}[Half-edge cancellation]
\[
\tau_{uv}+\tau_{vw}+\tau_{wu}=0.
\]
Each half-edge appears exactly twice across the triple overlap and cancels mod $2$.
\end{lemma}

\begin{definition}[Gauge normalization]
\[
\sigma_u^{\mathrm{norm}}(e):=\sigma_u(e)+f(u),
\]
where $f:B\to\mathbb F_2$ is chosen so that
\[
\sum_{e\ni u}\sigma_u^{\mathrm{norm}}(e)=0
\]
and fixed on a spanning tree.
\end{definition}

\begin{lemma}[Kernel triviality]
If $\Phi_{\mathrm{explicit}}([c])=0$, then $\sigma_u(e)=f(u)+f(v)$, hence $[c]=0$.
\end{lemma}

\begin{definition}[Cocycle lifting]
\[
\Psi([\beta])=[(\sigma_u^\beta)_u],\qquad
\sigma_u^\beta(e)=\beta(e)+f(u).
\]
\end{definition}

\begin{lemma}
\[
\Phi_{\mathrm{explicit}}\circ\Psi=\mathrm{id},\qquad
\Psi\circ\Phi_{\mathrm{explicit}}=\mathrm{id}.
\]
\end{lemma}

\begin{theorem}
\[
\mathcal H_R(G)\cong H^1(B;\mathbb F_2).
\]
\end{theorem}

\end{document}
